\section{Network classes and main results}\label{sec:model-results}

\subsection{Scalar leaky-recovery RFDE}

Let \(r=5\sqrt5\), \((\tau_0,\tau_1)=(4\sqrt5,5\sqrt5)\), and
\(H(v)=(v-1)^3\).  The autonomous scalar model is
\begin{align}
 \dot v={}&v-\frac{v^3}{3}-w
 +\eps\kappa_1\left\{\frac{v(t-\tau_0)+v(t-\tau_1)}2-v(t)\right\}
 \notag\\
 &+\eps\kappa_3\left\{
 \frac{H(v(t-\tau_0))+H(v(t-\tau_1))}{2}-H(v(t))\right\},\label{eq:scalar-v}\\
 \dot w={}&\eps(v-a-w).\label{eq:scalar-w}
\end{align}
The validated center is
\begin{equation}
 \eps=\frac15,\qquad a=\frac14,\qquad
 \kappa_1=\frac1{250},\qquad \kappa_3=\frac1{200}.       \label{eq:center}
\end{equation}
The full phase space is
\(X=C([-r,0],\R^2)\).  Because the recovery equation has no delayed recovery
slot, the positive-time semiflow factors through
\begin{equation}
 Y=C([-r,0],\R)\times\R,\qquad
 \Pi_Y(\phi_v,\phi_w)=(\phi_v,\phi_w(0)).               \label{eq:reduced-space}
\end{equation}
This reduction preserves every nonzero Floquet multiplier and pulls stable
sets back exactly; an obsolete recovery history is not an additional
crossing coordinate.

\subsection{Finite row-mass Dobrushin networks}

For a finite network let \(Q,B_0,B_1\in\R^{N\times N}\) be nonnegative
matrices and let \(\pi\) be any stationary probability vector satisfying
\begin{equation}
 Q\mathbf1=\mathbf1,\quad \pi^TQ=\pi^T,\quad
 \pi^T\mathbf1=1,\quad
 \tau(Q)\le\frac12,\qquad
 B_j\mathbf1=\frac12\mathbf1.                           \label{eq:row-mass-network}
\end{equation}
Here \(\tau(Q)\) is the Dobrushin coefficient.  Zero entries of \(\pi\)
are allowed, and no lower bound on \(\pi_{\min}\) enters the estimates.
The row-mass identities preserve the collective line
\(\Span\{\mathbf1\}\), but \(\Ker\pi^T\) need not be invariant under the
delayed layers: no identity \(\pi^TB_j=\frac12\pi^T\) is assumed.
Collective/quotient coordinates therefore produce an upper-triangular
operator.  The former balanced direct sum is the special case
\[
 \delta_B=0,\qquad
 \delta_B:=\delta_0+\delta_1,\qquad
 \delta_j:=\frac12\norm{\pi^TB_j-\tfrac12\pi^T}_1,
 \qquad 0\le\delta_B\le1.
\]

The finite leaky network to which the estimates below apply is, with
\(f(s)=s-s^3/3\) and both \(f\) and \(H\) acting componentwise,
\begin{align}
 \dot v(t)={}&f(v(t))-w(t)+3(Q-\Id)v(t)\notag\\
 &+\eps\kappa_1\{B_0v(t-\tau_0)+B_1v(t-\tau_1)-v(t)\}\notag\\
 &+\eps\kappa_3\{B_0H(v(t-\tau_0))+B_1H(v(t-\tau_1))-H(v(t))\},
                                                               \label{eq:leaky-network-v}\\
 \dot w(t)={}&\eps\{v(t)-a\mathbf1-w(t)\}+2(Q-\Id)w(t).
                                                               \label{eq:leaky-network-w}
\end{align}
Here \(v,w\in\R^N\), and all scalar constants are exactly those in
\cref{eq:center}.  The row-mass identities imply that
\(v=V\mathbf1\), \(w=W\mathbf1\) is invariant and that its restriction is
exactly \cref{eq:scalar-v,eq:scalar-w}, including the two half-weighted
delayed terms.  Thus the constants in the transverse theorem below refer to
this RFDE, not to an unspecified network realization.

Put \(P_\perp=\Id-\mathbf1\pi^T\) and
\[
 E_{N,\perp}=\Ker\pi^T\times\Ker\pi^T,
 \qquad
 m_N(x,y)=\max\{\operatorname{diam}x,3\operatorname{diam}y\}.
\]
For a complete synchronous trajectory \((V,W)\), the quotient variational
equation on \(E_{N,\perp}\) is
\begin{align}
 \dot x_\perp(t)={}&
 \{f'(V(t))-\eps\kappa_1-\eps\kappa_3H'(V(t))\}x_\perp(t)
 +3(Q-\Id)x_\perp(t)-y_\perp(t)\notag\\
 &+\eps\sum_{j=0}^1
 \{\kappa_1+\kappa_3H'(V(t-\tau_j))\}
 P_\perp B_jx_\perp(t-\tau_j),
                                                        \label{eq:transverse-variational}\\
 \dot y_\perp(t)={}&\eps x_\perp(t)
 +\{2(Q-\Id)-\eps\Id\}y_\perp(t).                    \label{eq:transverse-recovery}
\end{align}
The projection on the delayed terms is essential when
\(\pi^TB_j\ne\tfrac12\pi^T\).  Define
\(X_{N,1/10}=C([-r,0],E_{N,\perp})\) with the transverse history diameter
norm
\begin{equation}
 \norm{\phi}_{1/10}
 =\sup_{-r\le\theta\le0}e^{\theta/10}
   m_N(\phi_v(\theta),\phi_w(\theta)).                  \label{eq:diameter-norm}
\end{equation}
Let \(U_{\perp,N}(t,s)\) denote the evolution family of
\cref{eq:transverse-variational,eq:transverse-recovery} on this space.  For
\(F\in C_b(\R,E_{N,\perp})\), set
\(\norm F_{\perp,\infty}=\sup_t m_N(F_v(t),F_w(t))\).
Let \(L_{\perp,N}\) be the differential operator obtained by moving the
right-hand side of
\cref{eq:transverse-variational,eq:transverse-recovery} to the left, and set
\begin{align}
 \mathscr Y_{\perp,N}&=C_b(\R,E_{N,\perp}),\notag\\
 \mathscr X_{\perp,N}
 &=\{z\in C_b^1(\R,E_{N,\perp}):
          L_{\perp,N}z\in\mathscr Y_{\perp,N}\},
                                                               \label{eq:transverse-operator-spaces}\\
 \norm z_{\mathscr X_{\perp,N}}
 &=\sup_{t\in\R}\norm{z_t}_{1/10}
   +\norm{L_{\perp,N}z}_{\perp,\infty}.
\end{align}

\begin{theorem}[Complete-line transverse inverse]\label{thm:transverse}
Along every complete synchronous leaky trajectory satisfying
\(\sup_t\abs{V(t)-1}\le5/2\), every network satisfying
\cref{eq:row-mass-network} admits a unique bounded complete solution of the
forced transverse equation for every \(F\in C_b(\R,E_{N,\perp})\).  The
resulting causal Green operator
\[
 G_{\perp,N}:C_b(\R,E_{N,\perp})
 \longrightarrow\mathscr X_{\perp,N}
\]
and the evolution family satisfy
\begin{align}
 \norm{U_{\perp,N}(t,s)}_{1/10}
 &\le e^{-(t-s)/10},\qquad t\ge s,\notag\\
 \sup_t\norm{(G_{\perp,N}F)_t}_{1/10}
 &\le10\norm F_{\perp,\infty},\qquad
 \norm{G_{\perp,N}}_{\mathscr Y_{\perp,N}
      \to\mathscr X_{\perp,N}}\le11.                \label{eq:green-bound}
\end{align}
The constants are independent of the finite network size and admitted
topology.
\end{theorem}

The proof uses a directed Dobrushin--Halanay estimate and a forward pullback
of zero-history forced solutions.  It never inverts an RFDE semiflow on the
whole history space.

We separate the algebraic transfer from its leaky realization.  By
\cref{thm:transverse},
\(\cL_{\perp,N}:=L_{\perp,N}:\mathscr X_{\perp,N}
\to\mathscr Y_{\perp,N}\) is an isomorphism with inverse
\(G_{\perp,N}\); in the graph norm,
\(\norm{G_{\perp,N}}\le11\).

\begin{corollary}[Abstract upper-triangular root transfer]
\label{cor:canonical-transfer}
Let \(\mathscr X_\parallel,\mathscr Y_\parallel\) be Banach spaces, suppose
\(\cL_\parallel:\mathscr X_\parallel\to\mathscr Y_\parallel\) has
Fredholm index \(-1\) and one-dimensional cokernel spanned by
\(\psi\in\mathscr Y_\parallel^*\), normalized by
\(\psi(e_\parallel)=1\), and let
\(\Gamma_N\in\mathcal L(\mathscr X_{\perp,N},\mathscr Y_\parallel)\).
Define the bounded block operator
\[
 \cL_N:
 \mathscr X_\parallel\oplus\mathscr X_{\perp,N}
 \longrightarrow
 \mathscr Y_\parallel\oplus\mathscr Y_{\perp,N},
 \qquad
 \cL_N=
 \begin{pmatrix}\cL_{\parallel}&\Gamma_N\\0&\cL_{\perp,N}\end{pmatrix}.
\]
Let \(g_\parallel(\nu)\in\mathscr Y_\parallel\) define
\(d(\nu)=\psi(g_\parallel(\nu))\), and assume that \(d\) has a simple root
\(d(\nu_c)=0\).  Then
\begin{equation}
 T_N\cL_N=\cL_{\parallel}\oplus\cL_{\perp,N},           \label{eq:lin-sum}
\end{equation}
where
\[
 T_N(y_\parallel,y_\perp)
 =\bigl(y_\parallel-\Gamma_NG_{\perp,N}y_\perp,y_\perp\bigr).
\]
The full cokernel functional, normalized by
\(e_N=(e_\parallel,0)\), is
\begin{equation}
 \Psi_N(y_\parallel,y_\perp)
 =\psi\bigl(y_\parallel-\Gamma_NG_{\perp,N}y_\perp\bigr). \label{eq:full-cokernel}
\end{equation}
For the collective residual \((g_\parallel(\nu),0)\), one has
\(\Psi_N(g_\parallel(\nu),0)=d(\nu)\).  Hence
\(\cL_N\) has the same Fredholm index as \(\cL_\parallel\), while
\(d_N(\nu)=d(\nu)\) has the same root, slope, and orientation as the scalar
gap.
\end{corollary}

For a fixed scalar phase, endpoint, parameter, and gap realization, the leaky
network has precisely this block form, with all scalar data lifted
collectively and
\[
 \norm{\Gamma_NG_{\perp,N}}
 \le\frac{391}{2000}\delta_B.
\]
Thus \cref{thm:transverse} verifies the network block of the corollary.  The
result remains conditional on the scalar phase-fixed complete-history root
and the declared collective preparation.  The scalar leaky root itself and
an invariant asynchronous voltage strip are separate open problems.

\subsection{Shared-resource networks without a synchrony quotient}

We next state a separate network result whose node fields are heterogeneous
and whose slow recovery variable is shared.  For each \(N\ge2\), let
\(P_N\ge0\) be row stochastic and let \(\pi_N\) be a strictly positive
stationary probability vector.  For fixed \(D,\gamma>0\), assume
\begin{equation}
 P_N\mathbf1=\mathbf1,\qquad
 \pi_N^TP_N=\pi_N^T,\qquad \pi_N^T\mathbf1=1,\qquad
 \tau(P_N)\le1-\gamma.                                \label{eq:shared-markov}
\end{equation}
Set
\begin{equation}
 P_{c,N}=\mathbf1\pi_N^T,\qquad P_{\perp,N}=\Id-P_{c,N},\qquad
 E_N=\Ker\pi_N^T,\qquad A_N=D(P_N-\Id)|_{E_N},       \label{eq:shared-splitting}
\end{equation}
and equip \(\R^N\) with
\begin{equation}
 \norm{x}_N=\abs{\pi_N^Tx}+\osc(x).                   \label{eq:shared-norm}
\end{equation}
Let the curvature vectors satisfy, with constants independent of \(N\),
\begin{equation}
 0<c_-\le c_{i,N}\le c_+,\qquad \pi_N^Tc_N=\alpha>0. \label{eq:shared-curvature}
\end{equation}
Fix \(\beta>0\), \(K\ne0\), and finitely many delay locations
\(0\le\theta_0<\cdots<\theta_m\le\Theta_*\).  For
\begin{equation}
 B_{k,N}(\zeta)=B_{k,N}+\zeta R_{k,N},\qquad
 C_N(\zeta)=\sum_{k=0}^mB_{k,N}(\zeta),                \label{eq:shared-layers}
\end{equation}
assume the fixed-support operator bound and the atomwise full-row
neutrality
\begin{equation}
 \sup_N\sum_{k=0}^m
 \left(\norm{B_{k,N}}_{\cL(\R^N,\norm{\cdot}_N)}
       +\norm{R_{k,N}}_{\cL(\R^N,\norm{\cdot}_N)}\right)<\infty,
 \qquad \pi_N^TR_{k,N}=0.                             \label{eq:shared-layer-bounds}
\end{equation}
The latter is a row identity, not merely
\(\pi_N^TR_{k,N}\mathbf1=0\).

For \(\eps=\delta^2\), the RFDE is
\begin{align}
 \dot v(t)={}&\frac23\mathbf1-w(t)\mathbf1
 -c_N\circ(v(t)-\mathbf1)^{\circ2}
 -\frac\beta3(v(t)-\mathbf1)^{\circ3}+D(P_N-\Id)v(t)\notag\\
 &+\delta^2K\left\{C_N(\zeta)v(t)
 -\sum_{k=0}^mB_{k,N}(\zeta)
       v\left(t-\frac{\theta_k}{\delta}\right)\right\},
                                                               \label{eq:shared-rfde-v}\\
 \dot w(t)={}&\delta^2\{\pi_N^Tv(t)-a\},             \label{eq:shared-rfde-w}
\end{align}
where \(v\in\R^N\), \(w\in\R\), and the powers and products are
componentwise.  Its collective fold is
\((v,w,a)=(\mathbf1,2/3,1)\).  Write
\begin{equation}
 a=1+\delta^2\nu,\qquad \mu=a-1=\delta^2\nu,\qquad
 \dot M_{1,N}=\sum_{k=0}^m\theta_kR_{k,N},             \label{eq:shared-parameters}
\end{equation}
and define
\begin{align}
 g_N&=A_N^{-1}P_{\perp,N}c_N,\notag\\
 \kappa_N&=\frac\beta3+2\pi_N^T\diag(c_N)g_N,\qquad
 \nu_{0,N}=-\frac{3\kappa_N}{8\alpha^3},             \label{eq:shared-center}\\
 r_N&=\frac{K}{2\alpha}\pi_N^T\diag(c_N)
 A_N^{-1}P_{\perp,N}\dot M_{1,N}\mathbf1.            \label{eq:shared-rn}
\end{align}

We now fix the selection entering the theorem.  In the fold chart
\[
 s=\delta t,\qquad w=\frac23-\delta^2Y,\qquad
 v=\mathbf1+\delta\mathbf1X
       +\delta^2\{g_NX^2+h\},\qquad h\in E_N,
\]
the singular critical field is
\(q_0(X,Y)=(Y-\alpha X^2,-X)\), with singular orbit
\[
 \gamma_0(s)=\left(-\frac{s}{2\alpha},
                    \frac{s^2-2}{4\alpha}\right).
\]

\begin{definition}[Canonical preparation datum]\label{def:canonical-preparation}
Fix \(p>4\), sufficiently large for the finite tame losses below, and
\(B>2\Theta_*+2\).  A canonical preparation datum
\[
 \mathbf p_{\rm can}
 =(p,B,\chi_{\rm graph},\chi_{\rm plan},
   \sigma_h,\mathcal E_{\rm nor},h_*)
\]
consists of fixed graph-cutoff profiles, a planar joining cutoff, a
componentwise stable saturation, a degree-three normal extension operator,
and a common stable radius.  These objects and their required derivatives
through order twelve are fixed once and for all.

For each \(N\ge2\) and target \(0<\delta<1\), put
\[
 S_\delta=\sqrt{2p\log(1/\delta)},\qquad R_\delta=S_\delta+B.
\]
The datum generates a target-frozen preparation
\(\mathcal P^{\rm can}_{N,\delta}\).  It depends on \(N\) through the
physical coefficients, \(P_{\perp,N}\), \(A_N\), and \(g_N\), and on
\(\delta\) through \(S_\delta\) and \(R_\delta\), but its cutoff operators
are independent of \((\nu,\zeta)\).  Parameter derivatives are always taken
with the target-\(\delta\) preparation held fixed.

Let \(\mathcal U^{\rm pl}_{\rm ret}(\delta)\) be the shrinking planar
tubular neighborhood of
\(\gamma_0([-S_\delta-1,S_\delta+1])\), and let
\[
 \mathfrak H^{[2]}_\delta
 =\{\Phi_{q_0}^{-t}u:
       u\in\mathcal U^{\rm pl}_{\rm ret}(\delta),\
       0\le t\le2\Theta_*\}.
\]
The graph preparation agrees with the transformed physical RFDE on an open
fixed-buffer neighborhood of this continuous depth-two flow hull for every
current and delayed stable slot with \(\|h\|_N\le h_*\).  The planar
preparation agrees with the exact reduced graph field on the retained tube
and equals \(q_0\) near \(\gamma_0(\pm R_\delta)\).  The selected traces
satisfy
\[
 \mathscr H_\alpha(z^{a,\rm can}_N(-R_\delta))=0,
 \quad
 \mathscr H_\alpha(z^{r,\rm can}_N(R_\delta))=0,
 \quad
 X^{a,\rm can}_N(0)=X^{r,\rm can}_N(0)=0.
\]
\end{definition}

The normalized central gap of these traces is
\begin{align}
 \mathscr D^{\rm can}_N(\delta,\nu,\zeta)
  &=\frac{2}{\alpha e}\left\{
    \mathscr H_\alpha(z^{a,\rm can}_N(0))
   -\mathscr H_\alpha(z^{r,\rm can}_N(0))\right\},
                                                               \label{eq:shared-gap}\\
 \mathscr H_\alpha(X,Y)
  &=\frac12e^{-2\alpha Y}
    \left(\alpha Y-\alpha^2X^2+\frac12\right).       \label{eq:shared-first-integral}
\end{align}
On the retained graph, a zero of \(\mathscr D^{\rm can}_N\) is
equivalent to equality of the two selected complete RFDE histories.

\begin{theorem}[Main theorem: shared-resource network response]
\label{thm:shared-resource}
Assume \cref{eq:shared-markov,eq:shared-curvature,eq:shared-layer-bounds}
with all displayed bounds common in \(N\), and assume
\begin{equation}
 \inf_{N\ge2}\abs{r_N}>0.                              \label{eq:shared-nondegeneracy}
\end{equation}
Fix a compact interval \(I_\nu\) containing every \(\nu_{0,N}\) in its
interior with one common positive boundary distance, and one canonical
preparation datum \(\mathbf p_{\rm can}\) as in
\cref{def:canonical-preparation}.  Then there
exist \(\delta_0,\zeta_0,c_0,C>0\), depending only on the common model
bounds, that boundary distance, and \(\mathbf p_{\rm can}\), but independent
of \(N\), such that for every \(N\ge2\), every
\(0<\delta<\delta_0\), and every \(\abs\zeta<\zeta_0\), the gap
\(\mathscr D^{\rm can}_N/\delta\) has exactly one root in the local window
\(\{\nu:\abs{\nu-\nu_{0,N}}<c_0\}\subset I_\nu\),
\begin{equation}
 \nu^{\rm can}_{c,N}(\delta,\zeta)\in I_\nu,
 \qquad
 \abs{\nu^{\rm can}_{c,N}(\delta,\zeta)-\nu_{0,N}}<c_0. \label{eq:shared-root}
\end{equation}
No assertion is made here about additional zeros elsewhere in \(I_\nu\).
The selected attracting and repelling complete histories agree at this
root.  With
\(\mu^{\rm can}_{c,N}=\delta^2\nu^{\rm can}_{c,N}\),
\begin{align}
 &\left|\mu^{\rm can}_{c,N}(\delta,\zeta)
       -\mu^{\rm can}_{c,N}(\delta,0)
       -\mathscr C_N\delta^3\zeta\right|
 \le C\{\delta^4\abs\zeta+\delta^3\zeta^2\},        \label{eq:shared-response}\\
 &\mathscr C_N=-\frac{r_N}{\alpha}
 =-\frac{K}{2\alpha^2}\pi_N^T\diag(c_N)
 A_N^{-1}P_{\perp,N}\dot M_{1,N}\mathbf1.             \label{eq:shared-coeff}
\end{align}
In particular,
\(\inf_N\abs{\mathscr C_N}>0\).  The exact finite-\(\delta\) root is the
object selected by \(\mathbf p_{\rm can}\).  No uniform comparison with
other graph preparations is asserted.
\end{theorem}

The nondegeneracy class is nonempty without a hidden synchrony quotient.
Indeed, take \(\pi_N=N^{-1}\mathbf1\), \(P_N=P_{c,N}\), and
\[
 s_{i,N}=\frac{\sqrt3(2i-N-1)}{\sqrt{(N-1)(N+1)}},
 \qquad c_N=\mathbf1+\sigma s_N,\qquad 0<\sigma<1/\sqrt3.
\]
These choices satisfy
\[
 \pi_N^Ts_N=0,\qquad \pi_N^Ts_N^{\circ2}=1,\qquad
 \norm{s_N}_\infty<\sqrt3,\qquad \alpha=1.
\]
For two delays \(\theta_0<\theta_1\), set
\[
 B_{0,N}=B_{1,N}=bP_{c,N},\qquad
 R_{0,N}=s_N\pi_N^T,\qquad R_{1,N}=-s_N\pi_N^T,
 \qquad b>0.
\]
The layers remain entrywise positive for
\(\abs\zeta<b/\sqrt3\), all entries of
\(c_N\) are distinct, and
\begin{equation}
 \mathscr C_N=\frac{K\sigma(\theta_0-\theta_1)}{2D}\ne0
                                                               \label{eq:shared-witness-coeff}
\end{equation}
for every \(N\ge2\).  The distinct curvatures destroy every nontrivial
synchrony polydiagonal.

This instance proves that the general root mechanism is not merely an
arbitrary-size replication of a two-module quotient.  It does not cover
moving delay atoms, arbitrary perturbations of the base topology, or graph
families whose mixing gap closes, and it does not identify the
canonically selected root with an independently selected physical outer
canard.
